% !TEX root = ./信号与系统笔记.tex
% !TEX program = xelatex
%==================================================================
%  信号与系统笔记  ——  排版模板（本文件只放导言区，单独编译没有意义）
%  编译：xelatex 跑两遍（或 latexmk -xelatex）
%  依赖：TeX Live 2021+ 完整版；中文字体见下方 fonts 一节
%  上面两行 magic comment 是写给 VS Code / TeXstudio 的：
%  在这个文件里点“编译”会自动改去编译主文件 信号与系统笔记.tex，
%  否则编辑器会把本文件当根文件，报 “no output was made”。
%==================================================================
\documentclass[UTF8,11pt,oneside,openany,scheme=chinese,fontset=none]{ctexbook}

% ---------- 页面 ----------
\usepackage[a4paper,left=2.8cm,right=2.8cm,top=2.7cm,bottom=2.7cm,
            headheight=22pt,headsep=11pt,footskip=16pt]{geometry}

% ---------- 数学与图形 ----------
\usepackage{amsmath}
\usepackage{mathtools}
\usepackage{amsthm}
\usepackage{unicode-math}                 % OpenType 数学排版
\setmathfont{STIX Two Math}[Scale=1.09]   % 公式：STIX Two，笔画比 Times 重、对比强，整体放大近一档
\usepackage{tikz}
\usetikzlibrary{calc,positioning}
\usepackage{pgfplots}
\pgfplotsset{compat=1.18}
\usepackage{forest}
\usepackage{graphicx}
\usepackage{array,makecell,booktabs,tabularx,multirow}
\usepackage{enumitem}
\usepackage[most]{tcolorbox}
\usepackage{etoolbox}
\usepackage{fancyhdr}
\usepackage{titlesec}
\usepackage{titletoc}
\usepackage[font=small,labelfont={sf,bf},labelsep=quad,
            justification=centering,skip=6pt]{caption}
\usepackage{subcaption}                   % 并排子图 subfigure
\usepackage{float}                        % 需要“图就钉在这”时用 [H]
\usepackage[hidelinks,bookmarksnumbered]{hyperref}
% microtype 在 XeLaTeX 下会干扰中文断行与字距，故不启用
% \usepackage{microtype}

% ---------- 颜色 ----------
\definecolor{accent}{HTML}{2E4A63}   % 主色：深靛蓝（章节数字、线条）
\definecolor{warm}{HTML}{B4562C}     % 辅色：赤陶（提示、强调）
\definecolor{ink}{HTML}{1F2A33}      % 正文黑
\definecolor{muted}{HTML}{6B7B8A}    % 次要灰蓝
\definecolor{rule}{HTML}{D5DDE5}     % 细线
\definecolor{cool}{HTML}{F1F4F8}     % 定义盒底
\definecolor{cream}{HTML}{FBF5EC}    % 注解盒底
\definecolor{soft}{HTML}{F6F2EA}     % 代码底色
\definecolor{warnbg}{HTML}{FBEDE6}   % 注意盒底

% ---------- 字体 ----------
% 中文：霞鹜新致宋（正文）/ HarmonyOS Sans SC（标题）/ 霞鹜文楷（引文）
% 若本机没有这些字体，替换成「SimSun / SimHei / KaiTi」即可
\setmainfont{TeX Gyre Termes}              % 西文正文：Times 风格
\setsansfont{TeX Gyre Heros}[Scale=0.92]   % 西文无衬线
\setmonofont{Consolas}[Scale=0.85]
% 正文强调（\textbf）由伪粗宋体改为真·黑体：中文重音更醒目；斜体用楷体承担
\setCJKmainfont{LXGW Neo ZhiSong}[BoldFont={HarmonyOS Sans SC Bold},
                                  ItalicFont={LXGW WenKai GB Lite}]
\setCJKsansfont{HarmonyOS Sans SC}[AutoFakeBold=2.2]
\setCJKmonofont{HarmonyOS Sans SC}[AutoFakeBold=2.0]
\setCJKfamilyfont{zhkai}{LXGW WenKai GB Lite}[AutoFakeBold=1.6]
\newcommand{\kaishu}{\CJKfamily{zhkai}}

% ---------- 公式着色与编号 ----------
% 全部数学内容用主色深靛蓝，与墨色正文拉开层次
\AtBeginDocument{%
  \everymath{\color{accent}}%
  \everydisplay{\color{accent}}%
}
% 公式编号同步用主色
\newtagform{accent}[\color{accent}]{(}{)}
\usetagform{accent}

% ---------- 正文微调 ----------
\linespread{1.06}

% ---------- 常用记号 ----------
\newcommand{\dd}{\mathrm{d}}
\newcommand{\lrarr}{\;\longleftrightarrow\;}
\DeclareMathOperator{\Sa}{Sa}
\DeclareMathOperator{\Si}{Si}
\DeclareMathOperator{\sgn}{sgn}
\newcommand{\highlight}[1]{\colorbox{yellow!28}{\textbf{#1}}}
\newcommand{\ccr}[1]{\makecell{{\color{#1}\rule{1cm}{1cm}}}}   % 色卡（保留兼容）
\allowdisplaybreaks[2]

% ---------- 正文微调 ----------
\setlength{\parskip}{2pt plus 1pt}
\renewcommand{\arraystretch}{1.25}
\setlength{\abovedisplayskip}{7pt plus 2pt minus 2pt}
\setlength{\belowdisplayskip}{7pt plus 2pt minus 2pt}
\setlength{\abovedisplayshortskip}{4pt}
\setlength{\belowdisplayshortskip}{5pt}
\clubpenalty=1000 \widowpenalty=1000 \displaywidowpenalty=1000

%==================================================================
%  编（Part）：中文数字大号 + 编名 + 细线，独占一页
%==================================================================
\ctexset{part={name={第,编},number=\zhnum{part}}}
\titleformat{\part}[display]
  {\normalfont\centering}
  {\color{accent}\sffamily\bfseries\fontsize{54}{54}\selectfont 第\,\zhnum{part}\,编}
  {0.6em}
  {\sffamily\bfseries\color{ink}\fontsize{22}{28}\selectfont}
  [\vspace{0.6em}{\color{accent}\rule{3.4cm}{1.2pt}}]
\titlespacing*{\part}{0pt}{70pt}{50pt}

%==================================================================
%  章节标题
%==================================================================
% 一级：色块竖线 + 「第 N 章」小字 + 大标题 + 细线
\titleformat{\chapter}[display]
  {\normalfont\raggedright}
  {\vspace*{-10pt}{\color{accent}\rule[-2.5pt]{3pt}{13pt}}\hspace{7pt}%
   {\sffamily\bfseries\color{accent}\fontsize{11}{11}\selectfont 第\,\thechapter\,章}}
  {0.6em}
  {\sffamily\bfseries\color{ink}\fontsize{25}{31}\selectfont}
  [\vspace{0.7em}{\color{rule}\hrule height 0.7pt}]
\titlespacing*{\chapter}{0pt}{10pt}{26pt}

% 无编号的一级标题（目录、前言）
\titleformat{name=\chapter,numberless}[display]
  {\normalfont\raggedright}
  {\vspace*{-10pt}}
  {0pt}
  {\sffamily\bfseries\color{ink}\fontsize{24}{30}\selectfont}
  [\vspace{0.7em}{\color{rule}\hrule height 0.7pt}]
\titlespacing*{name=\chapter,numberless}{0pt}{10pt}{26pt}

% 二级：色块数字 + 无衬线标题 + 细线
\titleformat{\section}[block]
  {\normalfont\Large\sffamily\bfseries\color{accent}}
  {\raisebox{1.2pt}{\colorbox{accent}{\strut\hspace{2pt}\color{white}\thesection\hspace{2pt}}}}{0.7em}{}
  [\vspace{-0.15em}{\color{rule}\hrule height 0.7pt}\vspace{-0.1em}]
\titlespacing*{\section}{0pt}{1.5em}{0.9em}

% 三级：强调色编号
\titleformat{\subsection}[block]
  {\normalfont\large\sffamily\bfseries\color{ink}}
  {\color{warm}\thesubsection}{0.6em}{}
\titlespacing*{\subsection}{0pt}{1.3em}{0.6em}

% 四级：菱形符号
\titleformat{\subsubsection}[block]
  {\normalfont\normalsize\sffamily\bfseries\color{accent}}
  {\color{warm}\rule[-0.5pt]{2.2pt}{2.2pt}\hspace{2pt}\rule[-0.5pt]{2.2pt}{2.2pt}}{0.5em}{}
\titlespacing*{\subsubsection}{0pt}{1.1em}{0.5em}

%==================================================================
%  页眉页脚
%==================================================================
\renewcommand{\chaptermark}[1]{\markboth{\thechapter\enspace #1}{}}
\renewcommand{\sectionmark}[1]{\markright{\thesection\enspace #1}}

\pagestyle{fancy}
\fancyhf{}
\fancyhead[L]{\sffamily\footnotesize\color{muted}信号与系统笔记}
\fancyhead[R]{\sffamily\footnotesize\color{muted}\nouppercase{\leftmark}}
\fancyfoot[C]{\sffamily\footnotesize\color{muted}\thepage}
\renewcommand{\headrulewidth}{0.4pt}
\renewcommand{\headrule}{\hbox to\headwidth{\color{rule}\leaders\hrule height \headrulewidth\hfill}}
\fancypagestyle{plain}{%
  \fancyhf{}%
  \fancyfoot[C]{\sffamily\footnotesize\color{muted}\thepage}%
  \renewcommand{\headrulewidth}{0pt}}

%==================================================================
%  目录
%==================================================================
\renewcommand{\contentsname}{目录}
\titlecontents{part}[0em]
  {\addvspace{14pt}\sffamily\bfseries\large\color{accent}}
  {\contentslabel{5.0em}}
  {\hspace*{-5.0em}}
  {}

\titlecontents{chapter}[0em]
  {\addvspace{7.5pt}\sffamily\bfseries\color{ink}}
  {\contentslabel{5.6em}\color{accent}}
  {\hspace*{-5.6em}}
  {\hfill\hspace*{0.45em}\makebox[2.4em][r]{\color{accent}\thecontentspage}}
\titlecontents{section}[3.0em]
  {\addvspace{1.2pt}\small\color{ink}}
  {\contentslabel{2.4em}\color{muted}}
  {\hspace*{-2.4em}}
  {\titlerule*[0.6pc]{.}\hspace*{0.45em}\makebox[2.4em][r]{\color{muted}\thecontentspage}}
\titlecontents{subsection}[3.4em]
  {\addvspace{0.4pt}\small\color{muted}}
  {\contentslabel{3.2em}}
  {\hspace*{-3.2em}}
  {\titlerule*[0.6pc]{.}\hspace*{0.45em}\makebox[2.4em][r]{\color{muted}\thecontentspage}}

%==================================================================
%  每章导览小目录
%==================================================================
\newcommand{\chaptertoc}{%
  \startcontents[chapters]%
  \par\addvspace{6pt}\noindent%
  \begin{tcolorbox}[notebox,colback=cool,borderline west={2.4pt}{0pt}{accent},
                    left=13pt,right=13pt,top=7pt,bottom=7pt]%
    {\sffamily\bfseries\color{accent}本章导览}\par\vspace{1pt}%
    \printcontents[chapters]{}{1}{\setcounter{tocdepth}{2}}%
  \end{tcolorbox}%
}

%==================================================================
%  重点标注（颜色 + 加重）
%==================================================================
% \key  ：术语 / 结论关键词——主色加深加粗
% \point：易错点 / 例题标签——赤陶色加粗
% keyeq ：关键公式盒（无编号），keyeqn：关键公式盒（带编号）
\newcommand{\key}[1]{\textbf{\textcolor{accent}{#1}}}
\newcommand{\point}[1]{\textbf{\textcolor{warm}{#1}}}

\newenvironment{keyeq}{%
  \par\addvspace{8pt}\noindent%
  \begin{tcolorbox}[enhanced,breakable,boxrule=0pt,arc=1pt,
    colback=cool,borderline west={3pt}{0pt}{accent},
    left=16pt,right=14pt,top=7pt,bottom=7pt,ams equation*]%
}{\end{tcolorbox}\par\addvspace{8pt}}

\newenvironment{keyeqn}{%
  \par\addvspace{8pt}\noindent%
  \begin{tcolorbox}[enhanced,breakable,boxrule=0pt,arc=1pt,
    colback=cool,borderline west={3pt}{0pt}{accent},
    left=16pt,right=14pt,top=7pt,bottom=7pt,ams equation]%
}{\end{tcolorbox}\par\addvspace{8pt}}

%==================================================================
%  彩色盒子（定义 / 注解 / 注意 / 证明）
%==================================================================
\tcbset{
  notebox/.style={
    enhanced, breakable, boxrule=0pt, arc=2pt,
    left=11pt, right=11pt, top=8pt, bottom=8pt,
    fontupper=\normalsize, before skip=9pt, after skip=9pt,
    segmentation style={solid, rule, line width=0.5pt}
  }
}

% —— 定义（编号）——
\newcounter{definition}[chapter]
\newenvironment{definition}[1][]{%
  \refstepcounter{definition}\par\addvspace{2pt}%
  \begin{tcolorbox}[notebox,colback=cool,borderline west={2.4pt}{0pt}{accent}]%
  \noindent{\sffamily\bfseries\color{accent}定义\,\thechapter.\arabic{definition}}%
  \ifstrempty{#1}{}{\ \ {\sffamily\color{accent}（#1）}}%
  \par\nobreak\vspace{3pt}\nobreak\ignorespaces
}{%
  \end{tcolorbox}%
}

% —— 定义（不编号）——
\newenvironment{definition*}[1][]{%
  \par\addvspace{2pt}%
  \begin{tcolorbox}[notebox,colback=cool,borderline west={2.4pt}{0pt}{accent}]%
  \noindent{\sffamily\bfseries\color{accent}定义}%
  \ifstrempty{#1}{}{\ \ {\sffamily\color{accent}（#1）}}%
  \par\nobreak\vspace{3pt}\nobreak\ignorespaces
}{%
  \end{tcolorbox}%
}

% —— 注解 / 例题（米色）——
\newenvironment{note}[1][]{%
  \par\addvspace{2pt}%
  \begin{tcolorbox}[notebox,colback=cream,borderline west={2.4pt}{0pt}{warm},fontupper=\kaishu]%
  \ifstrempty{#1}{}{\noindent{\sffamily\bfseries\color{warm}#1}\par\vspace{3pt}}%
  \ignorespaces
}{%
  \end{tcolorbox}%
}

% —— 注意（赤陶）——
\newenvironment{warning}[1][]{%
  \par\addvspace{2pt}%
  \begin{tcolorbox}[notebox,colback=warnbg,borderline west={2.4pt}{0pt}{warm}]%
  \noindent{\sffamily\bfseries\color{warm}$\blacktriangleright$\,注意}%
  \ifstrempty{#1}{}{\ \ {\sffamily\color{warm}——#1}}%
  \par\nobreak\vspace{3pt}\nobreak\ignorespaces
}{%
  \end{tcolorbox}%
}

% —— 定理 / 例 / 备注（备后续章节使用）——
\newcounter{theorem}[chapter]
\newenvironment{theorem}[1][]{%
  \refstepcounter{theorem}\par\addvspace{2pt}%
  \begin{tcolorbox}[notebox,colback=white,borderline west={2.4pt}{0pt}{accent},
                    borderline north={0.5pt}{0pt}{rule},borderline south={0.5pt}{0pt}{rule}]%
  \noindent{\sffamily\bfseries\color{accent}定理\,\thechapter.\arabic{theorem}}%
  \ifstrempty{#1}{}{\ \ {\sffamily\color{accent}（#1）}}%
  \par\nobreak\vspace{3pt}\nobreak\ignorespaces
}{\end{tcolorbox}}

\newenvironment{example}[1][]{%
  \par\addvspace{2pt}%
  \begin{tcolorbox}[notebox,colback=white,borderline west={2.4pt}{0pt}{muted}]%
  \noindent{\sffamily\bfseries\color{muted}例}%
  \ifstrempty{#1}{}{\ \ {\sffamily\color{muted}#1}}%
  \par\nobreak\vspace{3pt}\nobreak\ignorespaces
}{\end{tcolorbox}}

% —— 证明（amsthm 定制：无缩进、小字距、醒目标签）——
\renewenvironment{proof}[1][\proofname]{%
  \par\addvspace{4pt}%
  \pushQED{\qed}%
  \normalfont\topsep2pt\trivlist
  \item[\hspace*{0pt}{\sffamily\bfseries\color{accent}#1}]\hspace*{0.4em}\ignorespaces
}{%
  \popQED\endtrivlist\@endpefalse
}

%==================================================================
%  代码（备用）
%==================================================================
\usepackage{listings}
\lstdefinestyle{plain}{
  backgroundcolor=\color{soft}, basicstyle=\ttfamily\small,
  keywordstyle=\color{accent}\bfseries, commentstyle=\color{muted}, stringstyle=\color{warm},
  frame=none, framesep=0pt, xleftmargin=22pt, xrightmargin=6pt,
  numbers=left, numberstyle=\ttfamily\scriptsize\color{muted}, numbersep=8pt,
  breaklines=true, showstringspaces=false, tabsize=2,
  aboveskip=9pt, belowskip=9pt,
  captionpos=t, extendedchars=false
}
\lstset{style=plain}
\renewcommand{\lstlistingname}{代码}

%==================================================================
%  封面信息
%==================================================================
\makeatletter
\newcommand{\subtitle}[1]{\gdef\@subtitle{#1}}
\gdef\@subtitle{}
\newcommand{\institute}[1]{\gdef\@institute{#1}}
\gdef\@institute{}
\newcommand{\cover}[1]{\gdef\@coverimage{#1}}
\gdef\@coverimage{cover.png}
\newcommand{\extrainfo}[1]{\gdef\@extrainfo{#1}}
\gdef\@extrainfo{}
\renewcommand{\maketitle}{%
  \begin{titlepage}%
    \begin{tikzpicture}[remember picture,overlay]
      % 顶部整幅插图（不裁切、不变形）
      \node[anchor=north west,inner sep=0pt] at (current page.north west)
        {\includegraphics[width=\paperwidth]{\@coverimage}};
      % 主标题区
      \node[anchor=north west,align=left,text width=15.4cm,inner sep=0pt]
        at ([xshift=2.8cm,yshift=-15.2cm]current page.north west) {%
          {\color{accent}\rule[-2pt]{2.6pt}{12pt}}\hspace{7pt}%
          {\sffamily\bfseries\color{accent}\fontsize{10.5}{10.5}\selectfont 学\ 习\ 笔\ 记}\par
          \vspace{16pt}
          {\sffamily\bfseries\color{ink}\fontsize{34}{40}\selectfont \@title}\par
          \vspace{10pt}
          {\sffamily\color{muted}\fontsize{13}{18}\selectfont \@subtitle}\par
          \vspace{26pt}
          {\color{accent}\rule{4.2cm}{1.6pt}}\par
          \vspace{14pt}
          {\sffamily\color{ink}\fontsize{11}{14}\selectfont \@author}\par
          \vspace{3pt}
          {\sffamily\color{muted}\fontsize{10}{13}\selectfont \@institute}\par
        };
      % 页脚下沿
      \node[anchor=south west,align=left,text width=15.4cm,inner sep=0pt]
        at ([xshift=2.8cm,yshift=3.0cm]current page.south west) {%
          {\color{rule}\rule{\linewidth}{0.5pt}}\par\vspace{6pt}
          {\sffamily\color{muted}\fontsize{9}{12}\selectfont \@date\hfill\@extrainfo}\par
        };
    \end{tikzpicture}%
  \end{titlepage}%
  \setcounter{page}{1}%
}
\makeatother

\graphicspath{{figures/}}

% ---------- 封面信息 ----------
\title{信号与系统笔记}
\subtitle{连续与离散系统的时域与频域分析}
\author{张竣翔}
\institute{中国传媒大学}
\date{\today}
\cover{cover.png}
\extrainfo{The starting point of the journey}
